%%%%%%%%%%%%%%%%%%%%%%%%%%%%%%%%%%%%%%%%%%%%%%%%%%%%%%%%%%%%%%%%%%%%%%%%%%%%%%
%
% Chapter file included in main project file using \input{}
%
% Assumes that LaTeX2e macros and packages defined in rgb_laser_physics.sty
%   are available
%
%%%%%%%%%%%%%%%%%%%%%%%%%%%%%%%%%%%%%%%%%%%%%%%%%%%%%%%%%%%%%%%%%%%%%%%%%%%%%%

 \chapter{One-dimensional Multi-Mode Laser Dynamics\label{chp:laser_dynamics_1d_mml}}

In this chapter, we describe the dynamics of one-dimensional laser amplifiers and oscillators by applying the quasi-normal mode expansions derived in \sct{laser_resonators_1d_qnm} to the wave equation given by \eqn{wave_eqn_1d} and the density matrix evolution equations defined by \eqn{fls_mbe_rwa_pol} and \eqn{fls_mbe_rwa_pop_diff}. Under the right experimental conditions, these multimode representations (approximate as they are) can provide remarkably illuminating descriptions of laser behavior, including optimum output coupling, frequency pulling, wave mixing, and mode-locking.

When multiple modes oscillate in a laser, they give rise to coherent modulations of the populations in the nonlinear gain medium that create interactions between those modes. The frequencies of these modulations are integer multiples of the free spectral range $\Delta \omega_\text{FSR}$ --- defined by \eqn{delta_w_fsr_def} --- between adjacent intracavity field modes. In \fig{multimode_gain_spectrum_1d}, we show a plot of a gain medium with a peak at frequency $\omega_0 = \omega_{a b}$ as a function of the frequency detuning. We have superimposed the frequency modal structure --- over several free-spectral ranges --- of a cavity containing that medium. In the following sections, we will find that for a particular frequency $\Delta \omega_q \equiv 2 q \pi$, fluctuations in the gain medium at frequency $2 (q - p) \pi$ couple the electric field amplitude with frequency $2 p \pi$ to the macroscopic polarization component at frequency $2 q \pi$.

 \begin{figure}
  \centering
  \includegraphics[width=4.5in]{figures/multimode_gain_spectrum_1d}
  \caption{\label{fig:multimode_gain_spectrum_1d} Plot of a gain medium with a peak at frequency $\omega_0 = \omega_{a b}$ as a function of the frequency detuning. We have superimposed the frequency modal structure --- over several free-spectral ranges --- of a cavity containing the medium. We will find that for a particular frequency $\Delta \omega_q \equiv 2 q \pi$, fluctuations in the gain medium at frequency $2 (q - p) \pi$ couple the electric field amplitude with frequency $2 p \pi$ to the macroscopic polarization component at frequency $2 q \pi$.}
 \end{figure}

% In \sct{laser_dynamics_1d_mml_qsl}, we will apply the \emph{strong} rate-equation approximation (REA) to construct a non-perturbative theory of a high-intensity pulsed laser that is driven by a short-duration pump and ``primed'' by a slowly-varying input field. In \sct{laser_dynamics_1d_mml_mll}, we will relax the REA to allow rapid intermodal interactions and build a weak-field perturbative model of \emph{coherent population pulsations}\index{Coherent population pulsations}\cite{ref:sargent1974lp} that lead to passive \emph{mode-locking}\index{Mode-locking} in either the time or the frequency domain.

\input{files/laser_dynamics_1d_mml_mbe}
\input{files/laser_dynamics_1d_mml_frq}
\input{files/laser_dynamics_1d_mml_qsl}
\input{files/laser_dynamics_1d_mml_mll}
\input{files/laser_dynamics_1d_mml_mfl}

%%%%%%%%%%%%%%%%%%%%%%%%%%%%%%%%%%%%%%%%%%%%%%%%%%%%%%%%%%%%%%%%%%%%%%%%%%%%%%%
%
% Chapter file included in main project file using \input{}
%
% Assumes that LaTeX2e macros and packages defined in rgb_laser_physics.sty
%   are available
%
%%%%%%%%%%%%%%%%%%%%%%%%%%%%%%%%%%%%%%%%%%%%%%%%%%%%%%%%%%%%%%%%%%%%%%%%%%%%%%

 \chapter{One-dimensional Multi-Mode Laser Dynamics\label{chp:laser_dynamics_1d_mml}}

In this chapter, we describe the dynamics of one-dimensional laser amplifiers and oscillators by applying the quasi-normal mode expansions derived in \sct{laser_resonators_1d_qnm} to the wave equation given by \eqn{wave_eqn_1d} and the density matrix evolution equations defined by \eqn{fls_mbe_rwa_pol} and \eqn{fls_mbe_rwa_pop_diff}. Under the right experimental conditions, these multimode representations (approximate as they are) can provide remarkably illuminating descriptions of laser behavior, including optimum output coupling, frequency pulling, wave mixing, and mode-locking.

When multiple modes oscillate in a laser, they give rise to coherent modulations of the populations in the nonlinear gain medium that create interactions between those modes. The frequencies of these modulations are integer multiples of the free spectral range $\Delta \omega_\text{FSR}$ --- defined by \eqn{delta_w_fsr_def} --- between adjacent intracavity field modes. In \fig{multimode_gain_spectrum_1d}, we show a plot of a gain medium with a peak at frequency $\omega_0 = \omega_{a b}$ as a function of the frequency detuning. We have superimposed the frequency modal structure --- over several free-spectral ranges --- of a cavity containing that medium. In the following sections, we will find that for a particular frequency $\Delta \omega_q \equiv 2 q \pi$, fluctuations in the gain medium at frequency $2 (q - p) \pi$ couple the electric field amplitude with frequency $2 p \pi$ to the macroscopic polarization component at frequency $2 q \pi$.

 \begin{figure}
  \centering
  \includegraphics[width=4.5in]{figures/multimode_gain_spectrum_1d}
  \caption{\label{fig:multimode_gain_spectrum_1d} Plot of a gain medium with a peak at frequency $\omega_0 = \omega_{a b}$ as a function of the frequency detuning. We have superimposed the frequency modal structure --- over several free-spectral ranges --- of a cavity containing the medium. We will find that for a particular frequency $\Delta \omega_q \equiv 2 q \pi$, fluctuations in the gain medium at frequency $2 (q - p) \pi$ couple the electric field amplitude with frequency $2 p \pi$ to the macroscopic polarization component at frequency $2 q \pi$.}
 \end{figure}

% In \sct{laser_dynamics_1d_mml_qsl}, we will apply the \emph{strong} rate-equation approximation (REA) to construct a non-perturbative theory of a high-intensity pulsed laser that is driven by a short-duration pump and ``primed'' by a slowly-varying input field. In \sct{laser_dynamics_1d_mml_mll}, we will relax the REA to allow rapid intermodal interactions and build a weak-field perturbative model of \emph{coherent population pulsations}\index{Coherent population pulsations}\cite{ref:sargent1974lp} that lead to passive \emph{mode-locking}\index{Mode-locking} in either the time or the frequency domain.

\input{files/laser_dynamics_1d_mml_mbe}
\input{files/laser_dynamics_1d_mml_frq}
\input{files/laser_dynamics_1d_mml_qsl}
\input{files/laser_dynamics_1d_mml_mll}
\input{files/laser_dynamics_1d_mml_mfl}

%%%%%%%%%%%%%%%%%%%%%%%%%%%%%%%%%%%%%%%%%%%%%%%%%%%%%%%%%%%%%%%%%%%%%%%%%%%%%%%
%
% Chapter file included in main project file using \input{}
%
% Assumes that LaTeX2e macros and packages defined in rgb_laser_physics.sty
%   are available
%
%%%%%%%%%%%%%%%%%%%%%%%%%%%%%%%%%%%%%%%%%%%%%%%%%%%%%%%%%%%%%%%%%%%%%%%%%%%%%%

 \chapter{One-dimensional Multi-Mode Laser Dynamics\label{chp:laser_dynamics_1d_mml}}

In this chapter, we describe the dynamics of one-dimensional laser amplifiers and oscillators by applying the quasi-normal mode expansions derived in \sct{laser_resonators_1d_qnm} to the wave equation given by \eqn{wave_eqn_1d} and the density matrix evolution equations defined by \eqn{fls_mbe_rwa_pol} and \eqn{fls_mbe_rwa_pop_diff}. Under the right experimental conditions, these multimode representations (approximate as they are) can provide remarkably illuminating descriptions of laser behavior, including optimum output coupling, frequency pulling, wave mixing, and mode-locking.

When multiple modes oscillate in a laser, they give rise to coherent modulations of the populations in the nonlinear gain medium that create interactions between those modes. The frequencies of these modulations are integer multiples of the free spectral range $\Delta \omega_\text{FSR}$ --- defined by \eqn{delta_w_fsr_def} --- between adjacent intracavity field modes. In \fig{multimode_gain_spectrum_1d}, we show a plot of a gain medium with a peak at frequency $\omega_0 = \omega_{a b}$ as a function of the frequency detuning. We have superimposed the frequency modal structure --- over several free-spectral ranges --- of a cavity containing that medium. In the following sections, we will find that for a particular frequency $\Delta \omega_q \equiv 2 q \pi$, fluctuations in the gain medium at frequency $2 (q - p) \pi$ couple the electric field amplitude with frequency $2 p \pi$ to the macroscopic polarization component at frequency $2 q \pi$.

 \begin{figure}
  \centering
  \includegraphics[width=4.5in]{figures/multimode_gain_spectrum_1d}
  \caption{\label{fig:multimode_gain_spectrum_1d} Plot of a gain medium with a peak at frequency $\omega_0 = \omega_{a b}$ as a function of the frequency detuning. We have superimposed the frequency modal structure --- over several free-spectral ranges --- of a cavity containing the medium. We will find that for a particular frequency $\Delta \omega_q \equiv 2 q \pi$, fluctuations in the gain medium at frequency $2 (q - p) \pi$ couple the electric field amplitude with frequency $2 p \pi$ to the macroscopic polarization component at frequency $2 q \pi$.}
 \end{figure}

% In \sct{laser_dynamics_1d_mml_qsl}, we will apply the \emph{strong} rate-equation approximation (REA) to construct a non-perturbative theory of a high-intensity pulsed laser that is driven by a short-duration pump and ``primed'' by a slowly-varying input field. In \sct{laser_dynamics_1d_mml_mll}, we will relax the REA to allow rapid intermodal interactions and build a weak-field perturbative model of \emph{coherent population pulsations}\index{Coherent population pulsations}\cite{ref:sargent1974lp} that lead to passive \emph{mode-locking}\index{Mode-locking} in either the time or the frequency domain.

\input{files/laser_dynamics_1d_mml_mbe}
\input{files/laser_dynamics_1d_mml_frq}
\input{files/laser_dynamics_1d_mml_qsl}
\input{files/laser_dynamics_1d_mml_mll}
\input{files/laser_dynamics_1d_mml_mfl}

%%%%%%%%%%%%%%%%%%%%%%%%%%%%%%%%%%%%%%%%%%%%%%%%%%%%%%%%%%%%%%%%%%%%%%%%%%%%%%%
%
% Chapter file included in main project file using \input{}
%
% Assumes that LaTeX2e macros and packages defined in rgb_laser_physics.sty
%   are available
%
%%%%%%%%%%%%%%%%%%%%%%%%%%%%%%%%%%%%%%%%%%%%%%%%%%%%%%%%%%%%%%%%%%%%%%%%%%%%%%

 \chapter{One-dimensional Multi-Mode Laser Dynamics\label{chp:laser_dynamics_1d_mml}}

In this chapter, we describe the dynamics of one-dimensional laser amplifiers and oscillators by applying the quasi-normal mode expansions derived in \sct{laser_resonators_1d_qnm} to the wave equation given by \eqn{wave_eqn_1d} and the density matrix evolution equations defined by \eqn{fls_mbe_rwa_pol} and \eqn{fls_mbe_rwa_pop_diff}. Under the right experimental conditions, these multimode representations (approximate as they are) can provide remarkably illuminating descriptions of laser behavior, including optimum output coupling, frequency pulling, wave mixing, and mode-locking.

When multiple modes oscillate in a laser, they give rise to coherent modulations of the populations in the nonlinear gain medium that create interactions between those modes. The frequencies of these modulations are integer multiples of the free spectral range $\Delta \omega_\text{FSR}$ --- defined by \eqn{delta_w_fsr_def} --- between adjacent intracavity field modes. In \fig{multimode_gain_spectrum_1d}, we show a plot of a gain medium with a peak at frequency $\omega_0 = \omega_{a b}$ as a function of the frequency detuning. We have superimposed the frequency modal structure --- over several free-spectral ranges --- of a cavity containing that medium. In the following sections, we will find that for a particular frequency $\Delta \omega_q \equiv 2 q \pi$, fluctuations in the gain medium at frequency $2 (q - p) \pi$ couple the electric field amplitude with frequency $2 p \pi$ to the macroscopic polarization component at frequency $2 q \pi$.

 \begin{figure}
  \centering
  \includegraphics[width=4.5in]{figures/multimode_gain_spectrum_1d}
  \caption{\label{fig:multimode_gain_spectrum_1d} Plot of a gain medium with a peak at frequency $\omega_0 = \omega_{a b}$ as a function of the frequency detuning. We have superimposed the frequency modal structure --- over several free-spectral ranges --- of a cavity containing the medium. We will find that for a particular frequency $\Delta \omega_q \equiv 2 q \pi$, fluctuations in the gain medium at frequency $2 (q - p) \pi$ couple the electric field amplitude with frequency $2 p \pi$ to the macroscopic polarization component at frequency $2 q \pi$.}
 \end{figure}

% In \sct{laser_dynamics_1d_mml_qsl}, we will apply the \emph{strong} rate-equation approximation (REA) to construct a non-perturbative theory of a high-intensity pulsed laser that is driven by a short-duration pump and ``primed'' by a slowly-varying input field. In \sct{laser_dynamics_1d_mml_mll}, we will relax the REA to allow rapid intermodal interactions and build a weak-field perturbative model of \emph{coherent population pulsations}\index{Coherent population pulsations}\cite{ref:sargent1974lp} that lead to passive \emph{mode-locking}\index{Mode-locking} in either the time or the frequency domain.

\input{files/laser_dynamics_1d_mml_mbe}
\input{files/laser_dynamics_1d_mml_frq}
\input{files/laser_dynamics_1d_mml_qsl}
\input{files/laser_dynamics_1d_mml_mll}
\input{files/laser_dynamics_1d_mml_mfl}

%\input{files/laser_dynamics_1d_mml}



